% Consolidated 08-relativity from 7c59eacf2bbc4f2718617404cd46e24a47721d7d; original TeX fragments preserved.

% === BEGIN 08-relativity-340-350.tex blob 8efb5bc3b2c2175c3e51a5b56863f6afa35c2cb3 ===
% Irodov, user-provided Arihant PDF, printed pp.56-57, PDF pp.65-66.
% Relativity Section 1.8. The original answer for #349 prints 2.2 x 10^4
% m/s, contradicting both its formula and the data; derived value ~2.2 x 10^8.
\exproblem{340}{56}
一根细杆沿自身长度方向以恒速 $v$
相对惯性参考系 $K$ 运动。
问速度为多少时，$K$ 系中测得的杆长
比其固有长度减少 $\eta=0.5\%$？
\exsolution
记杆固有长度 $L_0$，
运动方向的洛伦兹收缩为
\[
L=L_0\sqrt{1-\beta^2},\qquad\beta=v/c.
\]
长度比固有长度减少比例
$\eta=(L_0-L)/L_0$，故
$\sqrt{1-\beta^2}=1-\eta$。
两边平方：
\[
\beta^2=1-(1-\eta)^2=\eta(2-\eta).
\]
因此
\[
\boxed{v=c\sqrt{\eta(2-\eta)}
=c\sqrt{0.005(1.995)}
\simeq0.100c}.
\]
\excheck 相对速度虽然约为光速的 $10\%$，
长度收缩仍只有 $0.5\%$，因为低速收缩是
$\beta^2/2$ 量级。

\exproblem{341}{56}
一个等边三角形三边的固有长度都是 $a$。
在相对该三角形静止系以速率 $V$
作匀速运动的惯性系中，
若相对运动方向分别沿三角形
(a) 一条角平分线；(b) 一条边，
求三角形\emph{同一时刻测得}的周长；
并考察 $V\ll c$ 和 $V\to c$ 两种极限。
\exsolution
令 $\beta=V/c$，
沿相对运动方向的长度分量收缩因子
$q=\sqrt{1-\beta^2}$，垂直分量保持不变。
一条静止时与运动方向夹角为 $\theta$ 的边，
运动系长度为
\[
a'(\theta)=a\sqrt{q^2\cos^2\theta+\sin^2\theta}.
\]
(a) 角平分线方向与底边垂直，底边仍长 $a$；
两条腰与此方向分别成 $30^\circ$，
各腰测长为
$a\sqrt{(3/4)q^2+1/4}$。
于是
\[
\boxed{P_a=a\left[1+\sqrt{4-3\beta^2}\right]}.
\]
(b) 一条边平行于相对运动方向，
测长 $aq$；另两边各与运动方向成 $60^\circ$，
各测长 $a\sqrt{q^2/4+3/4}$，所以
\[
\boxed{P_b=a\left[
\sqrt{1-\beta^2}+\sqrt{4-\beta^2}\right]}.
\]
低速展开二者均为
$P=3a-\tfrac34a\beta^2+O(\beta^4)$。
当 $\beta\to1^-$ 时，
$\boxed{P_a\to2a,\quad P_b\to\sqrt3a}$。
\excheck 对每条边测量须在\emph{运动系同一时刻}
取两端坐标；不能直接把固有周长整体乘同一个
洛伦兹因子。

\exproblem{342}{56}
一根细杆在实验室系中以速度
$v=c/2$ 平动，其测得长度为
$L=1.00\,\mathrm m$；
杆与其速度方向的夹角
$\theta=45^\circ$。
求杆的固有长度 $L_0$。
\exsolution
在实验室中把杆长分解为
$L_\parallel=L\cos\theta$ 与
$L_\perp=L\sin\theta$。
洛伦兹收缩只发生在平行运动方向，
因此在杆静止系中
$L_{0,\parallel}=\gamma L_\parallel$，
$L_{0,\perp}=L_\perp$，
其中 $\gamma=1/\sqrt{1-\beta^2}$。
所以
\[
L_0^2=\gamma^2L^2\cos^2\theta+L^2\sin^2\theta
=L^2\frac{1-\beta^2\sin^2\theta}{1-\beta^2}.
\]
故
\[
\boxed{L_0=L\sqrt{
\frac{1-\beta^2\sin^2\theta}{1-\beta^2}}
=\sqrt{\frac{1-\frac14\cdot\frac12}{1-\frac14}}
\,\mathrm m\simeq1.08\,\mathrm m}.
\]
\excheck 原书答案排印中平方根结构可能不易辨认；
只有这一表达式满足长度量纲与平、垂直分量关系。

\exproblem{343}{56}
静止的直立圆锥顶角参数为
$\theta=45^\circ$（母线与轴线夹角），
侧面积 $S_0=4.0\,\mathrm{m^2}$。
在相对它以速度 $v=(4/5)c$
沿锥轴平移的惯性系中，求：
(a) 母线与轴线夹角 $\theta'$；
(b) 锥体侧面积 $S'$。
\exsolution
设锥底半径为 $R$、固有高度为 $H$，
于是 $\tan\theta=R/H$，
固有母线长 $s_0=\sqrt{R^2+H^2}$，
侧面积 $S_0=\pi R s_0$。
运动系中轴向高度收缩为
$H'=H\sqrt{1-\beta^2}$，
横向半径仍为 $R$。
因此
\[
\tan\theta'=\frac R{H'}
=\frac{\tan\theta}{\sqrt{1-\beta^2}},
\]
$\beta=4/5$、$\theta=45^\circ$ 时
$\tan\theta'=5/3$，故
$\boxed{\theta'\simeq59^\circ}$。
新的侧面积
$S'=\pi R\sqrt{R^2+H'^2}$，
以 $S_0$ 表示：
\[
\boxed{S'=S_0
\sqrt{1-\beta^2\cos^2\theta}
=4.0\sqrt{1-0.64(1/2)}
\simeq3.3\,\mathrm{m^2}}.
\]
\excheck 曲面面积不简单地统一乘以
$\sqrt{1-\beta^2}$；半径方向不收缩，
所以须重算母线长。

\exproblem{344}{56}
一只钟相对惯性系 $K$ 匀速运动。
$K$ 系中的标准钟经过
$t=5.0\,\mathrm s$ 时，
移动的钟只走过 $t-\Delta t$，
其中慢了 $\Delta t=0.10\,\mathrm s$。
求移动钟的速率 $v$。
\exsolution
移动钟所记录的是其世界线上的固有时间
$\tau=t-\Delta t$，时间膨胀关系为
\[
\tau=t\sqrt{1-v^2/c^2}.
\]
故
\[
\sqrt{1-\beta^2}=1-\frac{\Delta t}{t},
\quad
\beta=\sqrt{1-\left(1-\frac{\Delta t}{t}\right)^2}.
\]
代入 $\Delta t/t=0.10/5.0=0.020$，
\[
\boxed{v=c\sqrt{\frac{\Delta t}{t}
\left(2-\frac{\Delta t}{t}\right)}
\simeq0.199c
\simeq6.0\times10^7\,\mathrm{m/s}}.
\]
\excheck 必须比较同一对事件的运动钟固有时间和
$K$ 系坐标时间，不是“两只静止钟有不同走时率”。

\exproblem{345}{56}
一根杆沿自身长度方向相对于地面惯性系 $K$
以恒速飞过地面一个固定标记。
地面系中从杆头至杆尾通过标记历时
$\Delta t=20\,\mathrm{ns}$；
杆的静止系中，该标记通过整根杆历时
$\Delta t'=25\,\mathrm{ns}$。
求杆的固有长度 $L_0$。
\exsolution
记杆的速率为 $v$、洛伦兹因子 $\gamma$。
在 $K$ 中杆的长度收缩为 $L=L_0/\gamma$，
经过固定点所需时间
$\Delta t=L/v=L_0/(\gamma v)$。
在杆静止系中标记以速率 $v$
相对杆运动，走过固有长度所需
$\Delta t'=L_0/v$。
因此 $\gamma=\Delta t'/\Delta t=25/20=1.25$，
\[
\beta=\sqrt{1-\gamma^{-2}}
=\sqrt{1-\frac1{1.25^2}}=0.60.
\]
杆的固有长度
\[
\boxed{L_0=v\Delta t'
=(0.60c)(25\times10^{-9}\,\mathrm s)
\simeq4.5\,\mathrm m}.
\]
\excheck 两个时间间隔是在不同参考系测量的
“整根杆通过一个静止标记”的时间，
不可直接认为二者相等。

\exproblem{346}{56}
一颗不稳定粒子的固有寿命为
$\Delta\tau_0=10\,\mathrm{ns}$，
实验室系中从产生到衰变的时间间隔
$\Delta t=20\,\mathrm{ns}$。
求实验室中粒子在寿命内飞过的距离。
\exsolution
产生与衰变是粒子世界线上的两个事件，
粒子本身固有时间为 $\Delta\tau_0$。
相对实验室的洛伦兹因子
$\gamma=\Delta t/\Delta\tau_0=2$，
速度大小
\[
v=c\sqrt{1-\gamma^{-2}}
=c\frac{\sqrt3}{2}.
\]
实验室飞行距离为
\[
\boxed{L=v\Delta t
=c\Delta t\sqrt{
1-(\Delta\tau_0/\Delta t)^2}
\simeq5.2\,\mathrm m\approx5\,\mathrm m}.
\]
\excheck 固有寿命不是实验室中读出的寿命，
必须使用时间膨胀关系求粒子速度。

\exproblem{347}{56}
一颗 $\mu$ 子在实验室系中以
$v=0.990c$ 运动，
从产生点飞到衰变点的实验室距离为
$\ell=3.0\,\mathrm{km}$。
求：(a) $\mu$ 子的固有寿命；
(b) 在其自身静止系中，产生点与衰变点的相关
参照距离（等效为实验室中两固定位置的长度收缩值）。
\exsolution
在实验室中两事件相隔的坐标时间为
$\Delta t=\ell/v$。
粒子沿自身世界线经过的固有时间为
\[
\boxed{\Delta\tau_0=
\Delta t\sqrt{1-\beta^2}
=\frac{\ell}{v}\sqrt{1-\beta^2}}.
\]
代入 $\beta=0.990$，
$\sqrt{1-\beta^2}=\sqrt{0.0199}\simeq0.1411$，
得到 $\boxed{\Delta\tau_0\simeq1.43\,\mathrm{\mu s}}$。
在 $\mu$ 子系中，实验室的产生、衰变位置间
固有实验室尺距 $\ell$ 被收缩为
\[
\boxed{\ell'=\ell\sqrt{1-\beta^2}
\simeq0.423\,\mathrm{km}\approx0.42\,\mathrm{km}}.
\]
\excheck 在粒子自身系，粒子在产生与衰变之间
没有空间位移；上述 $\ell'$ 是\emph{移动的实验室两位置}
在粒子系中同一时刻测得的距离。
必须区分这一尺距与事件位移的概念。

\exproblem{348}{56}
两颗粒子在实验室中沿同一直线同向以
$v=(3/4)c$ 匀速运动，
先后击中同一个固定靶，两次击中
在实验室中相隔
$\Delta t=50\,\mathrm{ns}$。
求两粒子在其共同静止系中的固有间距。
\exsolution
两颗粒子在实验室中保持恒定间距
$L=v\Delta t$，因为同一靶被两者先后击中，
靶在两次碰撞间静止、相对粒子移动距离为 $v\Delta t$。
粒子共同静止系中的间距为固有长度，
实验室中沿运动方向收缩：
$L=L_0/\gamma$。
所以
\[
\boxed{L_0=\gamma v\Delta t
=\frac{v\Delta t}{\sqrt{1-v^2/c^2}}}.
\]
代入 $\beta=3/4$，
$\gamma=4/\sqrt7$，得到
$\boxed{L_0\simeq17\,\mathrm m}$。
\excheck 固有间距应比实验室中的
$v\Delta t\simeq11.25\,\mathrm m$ 更大。

\exproblem{349}{56}
一根杆沿一把静止尺匀速运动。
若杆两端坐标在\emph{尺的静止系同时}读取，
两刻度差为 $\Delta x_1=4.0\,\mathrm m$；
若杆两端坐标在\emph{杆自身静止系同时}读取，
读到的仍是尺刻度，但两位置对应的刻度差为
$\Delta x_2=9.0\,\mathrm m$。
求杆固有长度 $L_0$ 与相对尺的速率 $v$。
\exsolution
在尺静止系的同时测量给出长度收缩
$\Delta x_1=L_0/\gamma$。
在杆静止系的同时测量要求
$\Delta t'=0$，由逆洛伦兹变换，
这两个测量事件在尺系满足
$\Delta x=\gamma(\Delta x'+v\Delta t')
=\gamma L_0$，故
$\Delta x_2=\gamma L_0$。
联立得到
\[
\boxed{L_0=\sqrt{\Delta x_1\Delta x_2}
=\sqrt{4.0\times9.0}=6.0\,\mathrm m},
\]
\[
\gamma=\sqrt{\frac{\Delta x_2}{\Delta x_1}}
=\frac32,\qquad
\boxed{v=c\sqrt{1-\frac{\Delta x_1}{\Delta x_2}}
=c\sqrt{\frac59}\simeq2.24\times10^8\,\mathrm{m/s}}.
\]
\excheck \textbf{原书数值勘误：}本版答案页同时印有上述
正确公式和数值“$2.2\times10^4\,\mathrm{m/s}$”，
但 $c\sqrt{5/9}\simeq2.2\times10^8\,\mathrm{m/s}$，
相差 $10^4$ 倍；按公式和题面数据的结果已明确给出。

\exproblem{350}{57}
两根固有长度均为 $L_0$ 的杆沿共同轴线
彼此相向匀速运动。
在其中一根杆的静止系中，
两杆\emph{右端相遇}与\emph{左端相遇}
之间的时间间隔为 $\Delta t$。
求另一根杆相对于该杆的速率 $v$。
\exsolution
在第一杆静止系中，它长 $L_0$；
第二根杆沿轴运动且发生长度收缩，
长度为 $L=L_0\sqrt{1-\beta^2}$。
取两个端点的相遇顺序使移动杆在这段时间内
相对于第一杆扫过两杆长度之和，
因此
\[
v\Delta t=L_0+L
=L_0(1+\sqrt{1-\beta^2}).
\]
令 $a=L_0/(c\Delta t)$，
则 $\beta=a(1+\sqrt{1-\beta^2})$。
移项平方并剔除零速或不符合相遇顺序的分支，得
$\beta=2a/(1+a^2)$。
于是
\[
\boxed{v=\frac{2L_0/\Delta t}{
1+[L_0/(c\Delta t)]^2}}.
\]
\excheck 相遇方式必须是题定“右端相遇”和“左端相遇”
这两个不同事件；不能仅按一根杆的收缩长度
除以 $\Delta t$，会漏掉另一杆自身长度。
% === END 08-relativity-340-350.tex ===

% === BEGIN 08-relativity-351-360.tex blob b6ba04e71b2b87122109a492729be8f688799ee4 ===
% Irodov user PDF original pp.57-59, PDF pp.66-68.
% Special relativity: simultaneity, spacetime interval, velocity addition.
% Original diagrams 1.91-1.93 inspected.
\exproblem{351}{57}
在实验室惯性系 $K$ 中，两颗不稳定粒子沿
同一条直线同向以 $v=0.990c$ 运动，
任一相同实验室时刻它们的距离为 $\ell=120\,\mathrm m$。
若它们在\emph{共同静止系}中同时衰变，
求实验室系中两次衰变的时间差，
并说明哪颗粒子衰变得较晚。
\exsolution
取沿运动方向为正，前粒子比后粒子
在同一实验室时刻领先距离 $\ell$。
令前、后粒子衰变事件在 $K$ 系的
时空分离为 $\Delta t=t_f-t_r$、
$\Delta x=x_f(t_f)-x_r(t_r)$。
由于两粒子有相同速度，
$\Delta x=\ell+v\Delta t$。
共同静止系中的同时性要求
$\Delta t'=\gamma(\Delta t-v\Delta x/c^2)=0$。
代入：
\[
\Delta t=\frac v{c^2}(\ell+v\Delta t)
\Longrightarrow
\Delta t(1-v^2/c^2)=\frac{v\ell}{c^2}.
\]
故
\[
\boxed{\Delta t=\frac{\beta\ell}{c(1-\beta^2)}}
,\qquad\beta=v/c.
\]
代入 $\beta=0.990,\ell=120\,\mathrm m$：
\[
\boxed{\Delta t\simeq19.9\,\mathrm{\mu s}
\approx20\,\mathrm{\mu s}}.
\]
因为 $\Delta t>0$，实验室中\emph{前方粒子}
衰变得更晚。
\excheck 不可直接用 $\beta\ell/c$；
它忽略了同一实验室时刻的初始距离与两个衰变
事件发生在不同时刻之间的区别，漏掉 $1-\beta^2$。

\exproblem{352}{57}
杆 $AB$ 沿实验室 $K$ 系 $x$ 轴正方向以恒速 $v$
平动，前端为 $A$、后端为 $B$。
在时刻 $t_A$ 测得 $A$ 的坐标 $x_A$，
在时刻 $t_B$ 测得 $B$ 的坐标 $x_B$。
(a) 求杆固有长度 $L_0$；
(b) 若要求两个所记坐标\emph{之差的模}
恰等于固有长度，求可能的记录时间差。
\exsolution
记实验室中某一\emph{共同时间}的杆长为
$L=L_0\sqrt{1-\beta^2}$。
因 $A$ 匀速移动，
它在时刻 $t_B$ 的位置为
$x_A(t_B)=x_A-v(t_A-t_B)$。
故同时长度
\[
L=x_A-v(t_A-t_B)-x_B,
\]
因此
\[
\boxed{L_0=
\frac{x_A-x_B-v(t_A-t_B)}{\sqrt{1-\beta^2}}}.
\]
(b) 任意两个读数的有符号差满足
$x_A-x_B=L+v(t_A-t_B)$。
若要求 $x_A-x_B=+L_0$，则
\[
\boxed{t_A-t_B=\frac{L_0}{v}
\left(1-\sqrt{1-\beta^2}\right)}.
\]
若取相反顺序的差
$x_B-x_A=L_0$，则
\[
\boxed{t_B-t_A=\frac{L_0}{v}
\left(1+\sqrt{1-\beta^2}\right)}.
\]
\excheck 这两种记录时差正是原书要求的两个分支；
若只用同一时刻读数，则测得的是
$L=L_0/\gamma$，不可能等于 $L_0$（除非 $v=0$）。

\exproblem{353}{57}
两根固有长度均为 $L_0$ 的杆 $AB$
和 $A'B'$ 沿平行共线方向运动，后者相对前者
向右匀速运动 $v$（原书图1.91）。
每根杆两端各有彼此在\emph{该杆自身系内}同步的时钟。
规定 $B'$ 与 $A$ 相遇事件为两杆参考系时间原点。
求：(a) $B'$ 与 $B$ 相遇瞬间两只钟各自读数；
(b) $A'$ 与 $A$ 相遇瞬间两只钟各自读数。
\exsolution
以杆 $AB$ 静止系为 $K$，取 $A$ 在 $x=0$、
$B$ 在 $x=L_0$；令 $B'$ 与 $A$ 相遇为
$x=t=x'=t'=0$。
(a) $B'$ 在 $K$ 中以速度 $v$ 从 $A$ 走到 $B$
需时间 $t=L_0/v$。
静止钟 $B$ 读这个时间；随 $B'$ 运动的钟
读固有时间
\[
\boxed{t_B=\frac{L_0}{v}},\qquad
\boxed{t_{B'}=\frac{L_0}{\gamma v}
=\frac{L_0}{v}\sqrt{1-\beta^2}}.
\]
(b) 在 $K$ 中运动杆长度缩短为 $L_0/\gamma$，
因此 $A'$ 在初始实验室时刻位于
$-L_0/\gamma$；它追至静止的 $A$ 耗时
$t=L_0/(\gamma v)$。
故 $A$ 的静止钟读数为
\[
\boxed{t_A=\frac{L_0}{\gamma v}
=\frac{L_0}{v}\sqrt{1-\beta^2}}.
\]
该相遇事件的空间位置为 $x=0$，
洛伦兹时间变换给
$t'=\gamma(t-vx/c^2)=\gamma t=L_0/v$，
因此
\[
\boxed{t_{A'}=\frac{L_0}{v}}.
\]
\excheck 两对时钟读数的差别由时间膨胀和
\emph{同时性的相对性}共同造成，
不能直接假设两根杆上四只钟在 $K$ 系同时同步。

\exproblem{354}{57--58}
参考系 $K$、$K'$ 分别配有在各自系中同步的
一列空间时钟。$K'$ 相对 $K$ 沿 $+x$ 方向以
恒速 $v$ 平动（原书图1.92）。
在中央两只时钟 $A,A'$ 相遇且都读零的事件起，
分别画出：(a) 对 $K$ 系同时的时钟读数分布；
(b) 对 $K'$ 系同时的读数分布。
\exsolution
记 $\beta=v/c$、$\gamma=(1-\beta^2)^{-1/2}$，
两列时钟都选择中央相遇事件
$x=x'=t=t'=0$ 为零时刻。
(a) 在 $K$ 系的同时面 $t=0$ 上，
所有 $K$ 时钟同读 $0$。
$K'$ 时钟在本征坐标 $x'$ 处的读数由
洛伦兹变换
\[
t'=\gamma(t-vx/c^2),\qquad x'=\gamma x
\]
可得
\[
\boxed{t'(t=0,x')=-vx'/c^2}.
\]
因此沿 $K'$ 从左到右的钟依次为
\emph{超前、超前、零、落后、落后}。
(b) 在 $K'$ 同时面 $t'=0$ 上，
所有 $K'$ 钟同读 $0$。
由逆变换
$t=\gamma(t'+vx'/c^2)$、
$x=\gamma(x'+vt')$，
可得位于 $K$ 本征坐标 $x$ 的钟读数为
\[
\boxed{t(t'=0,x)=vx/c^2}.
\]
所以沿 $K$ 从左到右的钟依次为
\emph{落后、落后、零、超前、超前}。
为直观表达，以下为按图1.92重新绘制的概念时钟图，
左右两边指针的偏转只表示相对先后，不按数值比例：
\begin{center}
\begin{tikzpicture}[scale=0.82,>=Stealth]
\node[anchor=east] at (-3.3,0) {$K\text{ 同时}$};
\foreach \i in {-2,-1,0,1,2}{
  \pgfmathsetmacro{\xx}{1.3*\i}
  \pgfmathsetmacro{\ang}{90-20*\i}
  \draw (\xx,0) circle (0.34);
  \draw[->] (\xx,0) -- ++(\ang:0.28);
  \node[below] at (\xx,-0.45) {\scriptsize $t'$};
}
\node[anchor=east] at (-3.3,-2.0) {$K'\text{ 同时}$};
\foreach \i in {-2,-1,0,1,2}{
  \pgfmathsetmacro{\xx}{1.3*\i}
  \pgfmathsetmacro{\ang}{90+20*\i}
  \draw (\xx,-2) circle (0.34);
  \draw[->] (\xx,-2) -- ++(\ang:0.28);
  \node[below] at (\xx,-2.45) {\scriptsize $t$};
}
\end{tikzpicture}
\end{center}
\excheck 不能简单把两系全部时钟都画成零；
各自同步面不重合。上式的先后次序比示意指针
更精确且可定量重绘。

\exproblem{355}{58}
$K'$ 参考系相对 $K$ 沿 $+x$ 方向以恒速
$V$ 平动，两原点相遇时各自时钟读零。
在 $K$ 系中，找出一条匀速轨迹
$x=ut$，使它经过的每一对本地 $K$ 与 $K'$ 时钟
在相遇事件处的读数\emph{始终相等}。
求该轨迹速率 $u$，证明 $u<V$。
\exsolution
同一事件的坐标满足洛伦兹变换
$t'=\gamma(t-Vx/c^2)$。
要求在该事件读数相等即 $t'=t$，
所以
\[
\gamma(t-Vx/c^2)=t,
\qquad
x=\frac{c^2}{V}(1-\gamma^{-1})t.
\]
因此轨迹速率为
\[
\boxed{u=\frac{c^2}{V}
\left(1-\sqrt{1-V^2/c^2}\right)}.
\]
有理化得
\[
u=\frac{V}{1+\sqrt{1-V^2/c^2}}.
\]
因 $0<V<c$ 时分母严格大于 $1$，
\[
\boxed{0<u<V}.
\]
\excheck 这是两组同步时钟读数在\emph{对应相遇事件}
相等的轨迹，不是其中一组时钟的物理运动轨迹。

\exproblem{356}{58}
证明：若参考系 $K$ 中两个事件
存在真实因果联系，例如发射子弹与击中目标，
则在任何其他惯性系 $K'$ 中，
它们的发生先后顺序不能颠倒。
\exsolution
选两事件在 $K$ 系坐标差为
$\Delta t=t_2-t_1>0$、
$\Delta\bm x$。
因它们有因果联系，
任何连接它们的物理信号速率不能超过光速，
所以
$|\Delta\bm x|\le c\Delta t$。
令 $K'$ 相对 $K$ 以速度 $V$ 沿 $x$ 方向移动，
则
\[
\Delta t'=\gamma\left(
\Delta t-\frac{V\Delta x}{c^2}\right).
\]
利用 $|\Delta x|\le|\Delta\bm x|\le c\Delta t$：
\[
\Delta t'\ge
\gamma\Delta t\left(1-\frac{|V|}{c}\right)>0,
\]
因为惯性系间相对速率满足 $|V|<c$。
因此
\[
\boxed{t_2'>t_1'\quad\text{在任一惯性系成立}}.
\]
\excheck 仅对时间样或光样分离的因果关联可证明。
若事件为空间样分离，
不同惯性系可以给出不同先后顺序，
但它们不能彼此传递低于光速的因果信号。

\exproblem{357}{58}
原书图1.93的闵可夫斯基时空图中，
横轴 $x$、纵轴 $ct$ 均以米为单位。
三事件坐标为
$A:(2,1)$、$B:(5,6)$、$C:(7,4)$。
求：(a) 存在某惯性系使 $A,B$ 同地点发生时，
两事件的时间间隔；
(b) 存在另一惯性系使 $A,C$ 同时发生时，
两地点的空间间距。
\exsolution
时空间隔
$s^2=c^2(\Delta t)^2-(\Delta x)^2$
在洛伦兹变换下不变。
(a) $A\to B$：
$c\Delta t=6-1=5\,\mathrm m$，
$\Delta x=5-2=3\,\mathrm m$，
故 $s^2=25-9=16\,\mathrm{m^2}>0$，
是时间样分离。
在同地点参考系有 $\Delta x'=0$，
所以
\[
\boxed{\Delta t'=\frac{\sqrt{s^2}}{c}
=\frac{4\,\mathrm m}{c}
\simeq13.3\,\mathrm{ns}}.
\]
(b) $A\to C$：
$c\Delta t=4-1=3\,\mathrm m$，
$\Delta x=7-2=5\,\mathrm m$，
$s^2=9-25=-16\,\mathrm{m^2}<0$，
是空间样分离。在同时参考系
$\Delta t'=0$，于是
\[
\boxed{|\Delta x'|=\sqrt{-s^2}=4.0\,\mathrm m}.
\]
\excheck 本题两个子问的区分取决于
时空间隔的符号；读原图坐标不能把纵轴 $ct$
直接当作以秒为单位的时间。

\exproblem{358}{58}
惯性系 $K$ 中一粒子在 $xy$ 平面运动，
速度分量为 $v_x,v_y$。
$K'$ 相对 $K$ 沿 $+x$ 轴以速度 $V$ 平动，
两系坐标轴互相平行。
求粒子在 $K'$ 系中的速率 $v'$。
\exsolution
洛伦兹速度分量变换为
\[
v_x'=\frac{v_x-V}{1-v_xV/c^2},
\qquad
v_y'=\frac{v_y\sqrt{1-V^2/c^2}}{1-v_xV/c^2}.
\]
因为对有质量粒子
$|v_x|<c,|V|<c$，
分母 $1-v_xV/c^2>0$。
因此两分量平方相加并开方：
\[
\boxed{v'=\frac{
\sqrt{(v_x-V)^2+v_y^2(1-V^2/c^2)}
}{1-v_xV/c^2}}.
\]
\excheck 非平行方向速度不能简单用
$v'=|\bm v-\bm V|$，
横向分量也要经过洛伦兹因子和共同分母变换。

\exproblem{359}{58}
两粒子在实验室中沿同一直线相向运动，
各自速率为 $v_1=0.50c$ 与 $v_2=0.75c$。
求：(a) 实验室中两粒子间距缩小的速率；
(b) 一粒子在另一粒子静止系中测得的相对速率。
\exsolution
(a) 在实验室同时坐标中，
两粒子间距的变化率为两实验室速率之和：
\[
\boxed{v_{\rm closing}=v_1+v_2
=1.25c}.
\]
这是\emph{两物体间距的坐标变化率}，
不是某个物体相对一个惯性系的局部速度，
因此允许大于 $c$。
(b) 速度合成须用爱因斯坦速度加法：
\[
v_{\rm rel}=\frac{v_1+v_2}{1+v_1v_2/c^2}
=\frac{1.25c}{1+0.50(0.75)},
\]
即
\[
\boxed{v_{\rm rel}\simeq0.909c\approx0.91c}.
\]
\excheck 两种“速度”属于不同定义；
狭义相对论限制的是实际相对速度 $v_{\rm rel}<c$。

\exproblem{360}{58--59}
两根固有长度均为 $L_0$ 的杆，
在实验室中沿共同轴线以相同速率 $v$
彼此相向运动。
在其中一根杆的静止系中，
另一根杆的长度是多少？
\exsolution
取实验室中第一杆速度为 $+v$，
第二杆速度为 $-v$。
按一维速度叠加，
第二杆相对第一杆的速率
\[
w=\frac{(-v)-(+v)}{1-(-v)(+v)/c^2}
=-\frac{2v}{1+v^2/c^2}.
\]
令 $\beta=v/c$，相对速率的模为
$|w|/c=2\beta/(1+\beta^2)$，
对应洛伦兹因子
\[
\gamma_w=
\frac1{\sqrt{1-w^2/c^2}}
=\frac{1+\beta^2}{1-\beta^2}.
\]
因此另一根杆在第一根杆静止系中的收缩长度为
\[
\boxed{L=\frac{L_0}{\gamma_w}
=L_0\frac{1-\beta^2}{1+\beta^2}},
\qquad\beta=v/c.
\]
\excheck 不能直接套实验室速度 $v$ 的收缩因子，
因为两根杆间实际相对速率是
$2v/(1+\beta^2)$，而非 $2v$。
% === END 08-relativity-351-360.tex ===

% === BEGIN 08-relativity-361-370.tex blob ddc438496234589168f82ba01a9f3d65d3aeca7b ===
% Irodov, user-supplied Arihant PDF, original pp.59--60 (PDF pp.68--69).
% Section 1.8: special relativity. Problems 361--370, full derivations.
% Numerical constants: c=2.99792458e8 m/s, g=9.81 m/s^2.

\exproblem{361}{59}
两相对论性粒子在实验室系中沿相互垂直的方向运动，
速率分别为 $v_1,v_2$。
求两粒子之间的相对速率。
\exsolution
取粒子1的实验室速度沿 $+x$ 方向，即
$\bm v_1=v_1\bm e_x$；粒子2沿 $+y$，
$\bm v_2=v_2\bm e_y$。
转到粒子1的瞬时静止惯性系，其相对实验室系
沿 $+x$ 以 $v_1$ 运动。速度分量变换为
\[
v'_{2x}=\frac{0-v_1}{1-0}=-v_1,\qquad
v'_{2y}=\frac{v_2\sqrt{1-v_1^2/c^2}}{1-0}
=v_2\sqrt{1-v_1^2/c^2}.
\]
两分量平方求和，得两粒子的相对速率
\[
\boxed{v_{\rm rel}=
\sqrt{v_1^2+v_2^2-\frac{v_1^2v_2^2}{c^2}}}.
\]
\excheck 当任一速率为零时退化为另一速率，
而当二者均小于 $c$ 时，结果仍小于 $c$，
因为
$1-v_{\rm rel}^2/c^2=
(1-v_1^2/c^2)(1-v_2^2/c^2)>0$。

\exproblem{362}{59}
一个不稳定粒子在参考系 $K'$ 中沿 $y'$ 轴以
速率 $v'$ 运动。$K'$ 相对 $K$ 沿 $+x$ 轴
以速率 $V$ 匀速运动，$x,x'$ 轴同向、
$y,y'$ 轴平行。粒子的固有寿命为
$\Delta t_0$，求在 $K$ 系中粒子实际走过的距离。
\exsolution
令 $\beta=V/c$、
$\gamma_V=1/\sqrt{1-\beta^2}$、
$\gamma_{v'}=1/\sqrt{1-v'^2/c^2}$。
在 $K'$ 中粒子的速度为
$(v_x',v_y')=(0,v')$。
反向速度变换给
\[
v_x=V,\qquad v_y=\frac{v'}{\gamma_V},
\quad
v_K=\sqrt{V^2+(1-\beta^2)v'^2}.
\]
粒子在 $K'$ 中的寿命为
$\Delta t'=\gamma_{v'}\Delta t_0$。
世界线在 $K'$ 中满足 $\Delta x'=0$，
因此洛伦兹时间变换得到
$\Delta t=\gamma_V\Delta t'=
\gamma_V\gamma_{v'}\Delta t_0$。
于是在 $K$ 中的路程
\[
\boxed{s=v_K\Delta t=
\Delta t_0\sqrt{
\frac{V^2+(1-\beta^2)v'^2}
{(1-\beta^2)(1-v'^2/c^2)}}}.
\]
\excheck 当 $V=0$ 时
$s=v'\gamma_{v'}\Delta t_0$；
当 $v'=0$ 时
$s=V\gamma_V\Delta t_0$，
均符合单一参考系的时间膨胀关系。

\exproblem{363}{59}
一粒子在 $K$ 系内运动，速率为 $v$，
速度矢量与 $+x$ 轴夹角为 $\theta$。
$K'$ 系相对 $K$ 沿 $+x$ 方向
匀速运动，速率为 $V$，
两系 $x,x'$ 轴同向。
求在 $K'$ 系测得的运动方向角 $\theta'$。
\exsolution
在 $K$ 中
$v_x=v\cos\theta$、
$v_y=v\sin\theta$。
令 $\beta=V/c$，
洛伦兹速度变换为
\[
v_x'=\frac{v\cos\theta-V}
{1-vV\cos\theta/c^2},\qquad
v_y'=\frac{v\sin\theta\sqrt{1-\beta^2}}
{1-vV\cos\theta/c^2}.
\]
二者相除，得到
\[
\boxed{\tan\theta'=
\frac{\sqrt{1-\beta^2}\sin\theta}
{\cos\theta-V/v}}.
\]
\excheck 该式直接给出正切值；
若分母改变符号，必须依据
$(v_x',v_y')$ 所在象限选取
$\theta'$，不可只用 $\arctan$ 的主值。

\exproblem{364}{59}
如原书图1.94，一根细杆 $AB$ 在 $K'$ 系中
平行于 $x'$ 轴，且沿 $y'$ 负方向
以速度 $v'$ 平动。
与此同时，$K'$ 系相对实验室 $K$
沿 $+x$ 方向匀速运动，速率为 $V$。
求在 $K$ 系中\emph{同一时刻}看到的杆与
$x$ 轴之间的倾角大小 $\theta$。
\exsolution
杆在 $K'$ 的一个共同时间上两端满足
$\Delta x'=L'$、$\Delta y'=0$；
因杆还沿 $y'$ 方向平动，
两个任意取样事件满足
$\Delta y'=-v'\Delta t'$。
在实验室 $K$ 的同时性条件 $\Delta t=0$ 下，
洛伦兹变换
$\Delta t=\gamma_V(\Delta t'+V\Delta x'/c^2)$
给出
\[
\Delta t'=-\frac{V\Delta x'}{c^2},
\quad
\Delta y'=\frac{Vv'}{c^2}\Delta x'.
\]
实验室的两端横向距离
$\Delta y=\Delta y'$，
而
$\Delta x=\gamma_V(\Delta x'+V\Delta t')
=\Delta x'/\gamma_V$。
故
\[
\boxed{\tan\theta=
\gamma_V\frac{Vv'}{c^2}
=\frac{Vv'}{c^2\sqrt{1-V^2/c^2}}}.
\]
\excheck 这里的倾斜来自\emph{同时性相对性}：
实验室同时测量的两端在 $K'$ 中并非同时，
而杆在两次 $K'$ 时刻之间有横向位移。

\exproblem{365}{59}
惯性系 $K'$ 相对 $K$ 以恒速 $V$
运动。粒子在 $K$ 中沿一条直线运动，
瞬时速率为 $v$、加速度模为 $w$。
求它在 $K'$ 系的瞬时加速度 $w'$：
(a) 粒子速度平行于 $\bm V$；
(b) 粒子速度垂直于 $\bm V$。
\exsolution
取 $K'$ 沿 $+x$ 方向移动，
$\beta=V/c$，并从速度变换求导。
(a) 粒子沿 $x$ 运动，有
$v'_x=(v-V)/(1-vV/c^2)$。
对实验室时间求导：
\[
\frac{\mathrm dv'_x}{\mathrm dt}
=\frac{(1-\beta^2)w}{(1-vV/c^2)^2},
\quad
\frac{\mathrm dt'}{\mathrm dt}
=\gamma_V(1-vV/c^2).
\]
二者相除：
\[
\boxed{w'_\parallel=
\frac{w(1-\beta^2)^{3/2}}
{(1-vV/c^2)^3}}.
\]
若 $\bm v$ 与 $\bm V$ 反向，则
沿 $+x$ 的有符号 $v$ 应取负值。
(b) 若粒子沿 $y$ 运动并在该方向加速，
则实验室 $v_x=0$ 且始终 $a_x=0$。
$K'$ 中横向速度 $v'_y=v_y/\gamma_V$，
同时 $\mathrm dt'/\mathrm dt=\gamma_V$，
所以
\[
\boxed{w'_\perp=w(1-\beta^2)}.
\]
\excheck (a) 的因子取决于速度方向相对 $V$ 的
有符号关系；(b) 是粒子运动和加速度都垂直于
两系相对运动方向的情形。

\exproblem{366}{59}
一枚理想火箭从地球静止起飞，
其在每个瞬时共动惯性系中测得的
固有加速度均为 $w'=10g$，
推进阶段持续地球参考系时间
$\tau=1.0$ 年。
求推进末速比光速少百分之多少，以及
地球系中飞行的距离。
\exsolution
设固有加速度 $a=w'=10g$，
在地球系中火箭沿一直线加速。
一维加速度的变换关系为
$a=\gamma^3\,\mathrm dv/\mathrm dt$，于是
\[
\frac{\mathrm dv}{\mathrm dt}
=a\left(1-\frac{v^2}{c^2}\right)^{3/2}.
\]
由 $v(0)=0$ 积分：
\[
\frac{v}{\sqrt{1-v^2/c^2}}=at.
\]
解为
\[
\boxed{v(t)=\frac{at}{\sqrt{1+(at/c)^2}}}.
\]
因此末速相对光速的百分差为
\[
\boxed{100\frac{c-v(\tau)}c
=100\left[1-\frac{a\tau/c}{\sqrt{1+(a\tau/c)^2}}\right]
\simeq0.47\%}.
\]
对 $v(t)$ 再积分得
\[
\boxed{x(\tau)=\frac{c^2}{a}
\left[\sqrt{1+(a\tau/c)^2}-1\right]
\simeq0.91\ \text{光年}}.
\]
\excheck 一年在本题是\emph{地球系时间}，
不能拿它直接代入火箭钟的固有时长关系。
这里忽略引力及实际推进剂消耗。

\exproblem{367}{59}
沿用上一题的火箭与恒定固有加速度
$w'=10g$、地球时间 $\tau=1.0$ 年。
求推进阶段火箭自身时钟记录的固有时间
$\tau_0$。
\exsolution
由上一题速度规律
$v(t)=at/\sqrt{1+(at/c)^2}$，
得到
\[
\sqrt{1-v(t)^2/c^2}
=\frac1{\sqrt{1+(at/c)^2}}.
\]
按题给定义积分：
\[
\tau_0=\int_0^\tau
\sqrt{1-v(t)^2/c^2}\,\mathrm dt
=\int_0^\tau\frac{\mathrm dt}
{\sqrt{1+(at/c)^2}}.
\]
令 $u=at/c$，
利用 $\int\mathrm du/\sqrt{1+u^2}
=\operatorname{arsinh}u$，得
\[
\boxed{\tau_0=\frac ca
\operatorname{arsinh}\frac{a\tau}{c}
=\frac ca\ln\left[
\frac{a\tau}{c}+
\sqrt{1+\left(\frac{a\tau}{c}\right)^2}\right]}.
\]
代入 $a=10g$ 及地球年长，
$\boxed{\tau_0\simeq3.5\ \text{个月}}$。
\excheck 固有时间小于地球坐标时间；
积分不能把火箭的运动看作恒速时间膨胀。

\exproblem{368}{59}
一粒子速率比光速小 $0.010\%$，
求其所谓“相对论质量”
$m=\gamma m_0$ 与静质量 $m_0$ 之比。
\exsolution
速度亏缺分数
$\delta=(c-v)/c=0.010\%=10^{-4}$，
故 $\beta=v/c=1-\delta=0.9999$。
根据定义
\[
\frac{m}{m_0}=\gamma
=\frac1{\sqrt{1-\beta^2}}
=\frac1{\sqrt{1-(1-\delta)^2}}
=\frac1{\sqrt{\delta(2-\delta)}}.
\]
代入得
\[
\boxed{\frac{m}{m_0}\simeq70.7\approx70}.
\]
\excheck 现代物理通常优先采用\emph{不变静质量}
$m_0$，把 $\gamma$ 归于动量和能量，
而不是把惯性质量随速度增长当成普适基本量。
此处沿用原书术语并显式解释。

\exproblem{369}{60}
一物体在自身静止系中的质量密度为
$\rho_0$。
求在哪个相对该物体运动的惯性系中，
测得密度比 $\rho_0$ 大
$\eta=25\%$。
\exsolution
题目所指质量密度沿用原书
\emph{静质量除以同时测得的空间体积}的定义。
运动方向的线性尺寸收缩为
$L=L_0/\gamma$，横截面积不变，
故同一组静质量所占空间体积
$V=V_0/\gamma$。
于是
\[
\rho=\frac{m_0}{V}
=\gamma\frac{m_0}{V_0}=\gamma\rho_0.
\]
令 $\rho/\rho_0=1+\eta$，
$\gamma=1+\eta$，
\[
\boxed{v=c\sqrt{1-\frac1{(1+\eta)^2}}
=c\frac{\sqrt{\eta(2+\eta)}}{1+\eta}
=0.60c}.
\]
\excheck 这是\emph{静质量/收缩体积}的密度，
并不是包含运动动能的能量密度除以 $c^2$；
不同密度定义不可混用。

\exproblem{370}{60}
一颗质子的相对论动量为
$p=10.0\,\mathrm{GeV}/c$。
求其速率比光速小的百分比。
\exsolution
质子静止能量
$m_pc^2\simeq0.938\,\mathrm{GeV}$。
由质能关系
\[
E^2=p^2c^2+m_p^2c^4,\qquad v=\frac{pc^2}{E},
\]
得
\[
\beta=\frac vc
=\frac{pc}{\sqrt{(pc)^2+(m_pc^2)^2}}
=\left[1+\left(\frac{m_pc}{p}\right)^2\right]^{-1/2}.
\]
因此相对光速亏缺的百分比为
\[
\boxed{100(1-\beta)=100
\left[1-
\frac{1}{\sqrt{1+(0.938/10.0)^2}}\right]
\simeq0.44\%}.
\]
\excheck 高速下质子动量不能近似写作 $m_pv$；
此处用的是完整相对论能量—动量关系。
% === END 08-relativity-361-370.tex ===

% === BEGIN 08-relativity-371-380.tex blob 8a00871adc3a3ada7b7a17e223ddef5b2847a669 ===
% Irodov user-provided PDF pp.69--70, original printed pp.60--61.
% Relativity section 1.8, original problems 371--380.
\exproblem{371}{60}
质量为 $m_0$ 的粒子以何速率运动时，
其相对论动量模 $p$ 是牛顿动量模
$m_0v$ 的 $\eta=2.0$ 倍？
\exsolution
由相对论动量
$p=\gamma m_0v$，
其中 $\gamma=(1-\beta^2)^{-1/2}$、
$\beta=v/c$，
有
\[
\frac{p}{m_0v}=\gamma=\eta.
\]
因此
\[
\boxed{v=c\sqrt{1-\eta^{-2}}
=\frac c\eta\sqrt{\eta^2-1}
=\frac{\sqrt3}{2}c\simeq0.866c}.
\]
\excheck $p/(m_0v)$ 随速率增长且始终不小于1；
题给 $\eta=2$ 对应 $\gamma=2$，
与第373题的动能条件巧合相同。

\exproblem{372}{60}
把静质量为 $m_0$ 的粒子从
$v_1=0.60c$ 加速到 $v_2=0.80c$，
外力应做多少功？将相对论结果与
牛顿动能增量比较。
\exsolution
外力功等于总相对论能量的增量：
\[
W=(\gamma_2-\gamma_1)m_0c^2,
\quad
\gamma_i=\frac1{\sqrt{1-v_i^2/c^2}}.
\]
$\gamma_1=1/\sqrt{1-0.60^2}=5/4$，
$\gamma_2=1/\sqrt{1-0.80^2}=5/3$，
所以
\[
\boxed{W=\left(\frac53-\frac54\right)m_0c^2
=\frac5{12}m_0c^2\simeq0.42m_0c^2}.
\]
若误用牛顿动能，得到
\[
W_{\rm Newt}=\frac12m_0(v_2^2-v_1^2)
=\tfrac12m_0c^2(0.64-0.36)
=\boxed{0.14m_0c^2}.
\]
\excheck 牛顿估计仅约为正确值的三分之一，
说明在 $0.6c$ 到 $0.8c$ 的速度范围
不能忽略相对论效应。

\exproblem{373}{60}
粒子速率达到多少时，其动能
$T$ 恰好等于其静止能量 $E_0=m_0c^2$？
\exsolution
粒子的动能为
$T=(\gamma-1)m_0c^2$。
令 $T=m_0c^2$，得 $\gamma=2$。
由
$\gamma^{-2}=1-\beta^2$，
\[
\boxed{v=c\sqrt{1-\frac1{\gamma^2}}
=\frac{\sqrt3}{2}c
\simeq2.60\times10^8\,\mathrm{m/s}}.
\]
\excheck 动能等于静能并不表示粒子达到光速；
有静质量粒子的速率始终小于 $c$。

\exproblem{374}{60}
当粒子动能与静止能量之比
$q=T/(m_0c^2)$ 足够小时，
能否用牛顿公式 $T=m_0v^2/2$ 估算其速度？
求使速率相对误差小于
$\varepsilon=0.010$ 的 $q$ 上界近似值。
\exsolution
精确相对论关系
$T=(\gamma-1)m_0c^2$
给出 $\gamma=1+q$，从而
\[
v=c\sqrt{1-(1+q)^{-2}}.
\]
由牛顿动能反算的速率为
$v_{\rm N}=c\sqrt{2q}$。
对 $q\ll1$，
\[
v=c\sqrt{2q}
\left(1-\frac34q+O(q^2)\right).
\]
于是采用真实速度作相对误差分母，
\[
\frac{v_{\rm N}-v}{v}
=\frac34q+O(q^2).
\]
若要求误差小于 $\varepsilon$，
低能极限需要
\[
\boxed{\frac{T}{m_0c^2}=q
\lesssim\frac43\varepsilon
=\frac43(0.010)
\simeq0.013}.
\]
为避免将一阶估计误当成严格上界，
也可求\emph{精确}允许范围。利用
\[
\frac{v_{\rm N}}{v}
=\sqrt{\frac{2q}{1-(1+q)^{-2}}}
=(1+q)\sqrt{\frac{2}{2+q}},
\]
严格条件是
\[
\boxed{\frac{2(1+q)^2}{2+q}<(1+\varepsilon)^2}.
\]
取$\varepsilon=0.010$，该不等式对应
\[
\boxed{0<q<q_* \simeq0.0133704},\qquad
q_*=\frac{(\varepsilon-1)(\varepsilon+3)
 +(1+\varepsilon)\sqrt{\varepsilon^2+2\varepsilon+9}}4.
\]
\excheck 原书所用$q\lesssim 0.013$是
低能小参数的一阶近似；实际严格阈值为
$q_*\simeq0.0133704$。
须声明所用相对误差分母是\emph{真实相对论速度}，
若改用牛顿速度作分母，阈值会不同。

\exproblem{375}{60}
推导静质量为 $m_0$ 的粒子的
相对论动量 $p$ 与动能 $T$ 的关系。
并计算动能为 $500\,\mathrm{MeV}$
的质子的动量模。
\exsolution
相对论总能量
$E=T+m_0c^2$，
满足
$E^2=p^2c^2+m_0^2c^4$。
故
\[
p^2c^2=(T+m_0c^2)^2-m_0^2c^4
=T(T+2m_0c^2),
\]
因此
\[
\boxed{p=\frac1c
\sqrt{T(T+2m_0c^2)}}.
\]
质子静止能量
$m_pc^2\simeq938\,\mathrm{MeV}$，
代入 $T=500\,\mathrm{MeV}$：
\[
pc=\sqrt{500(500+2\cdot938)}\,\mathrm{MeV}
\simeq1090\,\mathrm{MeV},
\]
\[
\boxed{p\simeq1.09\,\mathrm{GeV}/c}.
\]
\excheck 若使用牛顿公式 $p=\sqrt{2m_0T}$
会漏掉 $T^2$ 项，无法在高动能下保证精度。

\exproblem{376}{60}
一束动能为 $T$ 的相对论性带电粒子
垂直射向完全吸收靶，
束流电流为 $I$，
每粒子电荷大小为 $e$、静质量为 $m_0$。
求靶受到的平均冲击力，以及靶释放的功率。
若给定受照面积 $A$，求平均压强。
\exsolution
单位时间打到靶上的粒子数为
$\dot N=I/e$。
粒子被完全吸收时，其初始动量全部传给靶，
每粒子动量模由能量—动量关系为
\[
p=\frac1c\sqrt{T(T+2m_0c^2)}.
\]
平均冲击力是动量通量：
\[
\boxed{F=\frac{I}{ec}
\sqrt{T(T+2m_0c^2)}}.
\]
若束斑面积为 $A$，
压力（压强）为
$\boxed{P_{\rm press}=F/A}$；
\emph{原题未提供面积，不能单独算出数值压强}。
假定打到靶上的粒子仅把其\emph{动能}
转化为热，释放的功率
\[
\boxed{\mathcal P=\dot NT=\frac{IT}{e}}.
\]
\excheck 原书简答使用符号 $F$ 表示作用力，
而不是面积归一化压强。若讨论粒子静质量
经反应转换为热，还需给出具体反应过程，
本题只考虑吸收、停止。

\exproblem{377}{60}
一球以相对论速率 $v$ 在稀薄气体中运动，
气体静止系每单位体积含 $n$ 个慢速粒子，
各粒子静质量为 $m$。
假设气体粒子在球表面垂直于球速的微面元上
发生理想弹性镜面反射。
求该面元受到的气体压强 $p$，
并验证在球的瞬时静止系和气体静止系所得相同。
\exsolution
记 $\gamma=(1-v^2/c^2)^{-1/2}$。
先到\emph{球的静止系}：
气体朝向表面的速度为 $v$，
长度收缩使粒子数密度变为
$n'=\gamma n$。
每粒子入射动量模为
$p_i=\gamma mv$，弹性反射使法向动量反号，
故每次冲量为 $2\gamma mv$。
单位面积单位时间碰撞数为
$n'v=\gamma nv$，因此
\[
\boxed{p=(\gamma nv)(2\gamma mv)
=\frac{2nmv^2}{1-v^2/c^2}}.
\]
在\emph{气体静止系}中，靶面以速度 $v$
迎向静止粒子。反射后粒子的速率大小由
相对论速度合成得
$u=2v/(1+v^2/c^2)$，
其动量为
$\gamma_u m u=2\gamma^2mv$。
单位面积单位时间被迎面扫过的粒子数为 $nv$，
所以靶面所得冲量通量仍为
\[
(nv)(2\gamma^2mv)=2\gamma^2nmv^2.
\]
\excheck 这里是法向的\emph{压力即动量流率密度}；
在选定的理想单向弹性反射模型下两惯性系计算一致。
粒子不能被当作“用牛顿动量 $mu$”来计算，
否则少一个相对论因子。

\exproblem{378}{60}
静质量为 $m_0$ 的粒子从
$t=0$ 时刻由静止开始受到
沿直线方向大小恒定的力 $F$。
求粒子速率 $v(t)$ 及在力方向的位移 $x(t)$。
\exsolution
沿运动方向的相对论动量为
$p=\gamma m_0v$，
动力学方程 $\mathrm dp/\mathrm dt=F$。
由 $p(0)=0$ 积分：
$p(t)=Ft$。
又
$v=pc^2/\sqrt{m_0^2c^4+p^2c^2}$，
因此
\[
\boxed{v(t)=
\frac{F t}{\sqrt{m_0^2+(F t/c)^2}}}.
\]
从 $x(0)=0$ 出发，
\[
x(t)=\int_0^t\frac{F s\,\mathrm ds}
{\sqrt{m_0^2+(Fs/c)^2}}
=\frac{m_0c^2}{F}
\left[\sqrt{1+\left(\frac{F t}{m_0c}\right)^2}-1\right].
\]
故
\[
\boxed{x(t)=\frac{
\sqrt{m_0^2c^4+F^2c^2t^2}-m_0c^2}{F}}.
\]
\excheck 在 $Ft\ll m_0c$ 时，
$v\simeq(F/m_0)t$，
$x\simeq Ft^2/(2m_0)$；
而 $t\to\infty$ 时
$v\to c^-$，始终不超光速。

\exproblem{379}{60}
一个粒子在 $K$ 系的运动规律为
$x(t)=\sqrt{a^2+c^2t^2}$，
其中 $a>0$ 为常长度、
$c$ 是光速。
求作用于它的相对论力 $F$。
\exsolution
直接微分：
\[
v=\dot x=\frac{c^2t}{x},\qquad
w=\ddot x=\frac{c^2a^2}{x^3}.
\]
由
$1-v^2/c^2=1-c^2t^2/x^2=a^2/x^2$
得洛伦兹因子
$\gamma=x/a$。
沿速度方向的动量微分
\[
F=\frac{\mathrm d}{\mathrm dt}(\gamma m_0v)
=m_0\gamma^3 w
=m_0\left(\frac xa\right)^3
\frac{c^2a^2}{x^3}.
\]
故作用力为恒定值
\[
\boxed{F=\frac{m_0c^2}{a}}
\]
且沿 $+x$ 方向。
\excheck 虽然坐标加速度 $w(t)$ 不恒定，
相对论力 $F=\mathrm dp/\mathrm dt$ 却恒定；
不能将牛顿式 $F=m_0w$ 直接使用。

\exproblem{380}{60--61}
从相对论动力学基本方程出发：
(a) 说明哪些情况下粒子的加速度方向
与作用力 $\bm F$ 相同；
(b) 当 $\bm F\perp\bm v$ 与
$\bm F\parallel\bm v$ 时，
求力与加速度模之间的比例因子。
\exsolution
相对论动量
$\bm p=\gamma m_0\bm v$。
求导并使用
$\dot\gamma=\gamma^3(\bm v\cdot\bm w)/c^2$，
得到
\[
\bm F
=\frac{\mathrm d\bm p}{\mathrm dt}
=\gamma m_0\bm w+
\frac{\gamma^3m_0}{c^2}
(\bm v\cdot\bm w)\bm v.
\]
将加速度分解为沿速率方向和法向分量
$\bm w=\bm w_\parallel+\bm w_\perp$，
可化为
\[
\bm F_\perp=\gamma m_0\bm w_\perp,\qquad
\bm F_\parallel=\gamma^3m_0\bm w_\parallel.
\]
由于两个系数通常不同，
$\bm F$ 与 $\bm w$ 在一般斜向情形并不同向。
仅在 $\bm F\parallel\bm v$
或 $\bm F\perp\bm v$
（以及 $v=0$ 的退化瞬间）时，二者同向。
因此
\[
\boxed{
F_\perp=\frac{m_0w_\perp}{\sqrt{1-\beta^2}},
\qquad
F_\parallel=\frac{m_0w_\parallel}{(1-\beta^2)^{3/2}}
},\qquad \beta=v/c.
\]
\excheck 纵向“等效惯性系数”为 $\gamma^3m_0$，
横向为 $\gamma m_0$；这些是方向依赖的动力学关系，
不是两个不同的不变静质量。
% === END 08-relativity-371-380.tex ===

% === BEGIN 08-relativity-381-388.tex blob 5da2171cabb568878d09975fabed50cc4a4de06a ===
% User-provided Irodov PDF pp.70, printed p.61;
% Section 1.8 final problems 381--388.
% Source answer p.295 has a unit error in problem 386: 1.43e3 GeV,
% not 1.43e3 MeV, for T=25 GeV per proton.
\exproblem{381}{61}
一个相对论性粒子在惯性系 $K$ 中沿
$x$ 轴运动，动量为 $p_x$、总能量为 $E$。
证明在相对 $K$ 沿 $+x$ 方向以恒速 $V$
运动的另一惯性系 $K'$ 中，
粒子动量和总能量变换为
\[
p'_x=\frac{p_x-EV/c^2}{\sqrt{1-V^2/c^2}},
\qquad
E'=\frac{E-Vp_x}{\sqrt{1-V^2/c^2}}.
\]
\exsolution
粒子在自身世界线上的固有时间
$\mathrm d\tau$ 为洛伦兹标量。
四动量定义为
\[
P^\mu=m_0\frac{\mathrm dx^\mu}{\mathrm d\tau},
\quad
x^\mu=(ct,x,y,z),
\quad
P^\mu=(E/c,p_x,p_y,p_z).
\]
坐标的洛伦兹变换为
\[
ct'=\gamma_V(ct-\beta x),\qquad
x'=\gamma_V(x-Vt),
\quad
\gamma_V=(1-\beta^2)^{-1/2},\quad\beta=V/c.
\]
由于分母 $\mathrm d\tau$ 不变，
对坐标变换沿世界线求导可知四动量
具有完全相同的线性变换律：
\[
P'^0=\gamma_V(P^0-\beta P^1),\qquad
P'^1=\gamma_V(P^1-\beta P^0).
\]
代回物理分量：
\[
\boxed{E'=\gamma_V(E-Vp_x)},\qquad
\boxed{p'_x=\gamma_V(p_x-VE/c^2)}.
\]
\excheck 四动量变换同时适用于有质量粒子和光子；
不能将能量视为标量而使其在各惯性系不变。

\exproblem{382}{61}
一光子在惯性系 $K$ 中的能量为
$\varepsilon$。
惯性系 $K'$ 沿光子传播方向以速度 $V$
相对 $K$ 运动。
用上一题四动量变换求光子在 $K'$ 中的能量
$\varepsilon'$；并求使
$\varepsilon'=\varepsilon/2$ 的 $V$。
\exsolution
光子沿正向传播，动量
$p_x=\varepsilon/c$。
由四动量能量变换，
\[
\varepsilon'=\gamma_V(\varepsilon-Vp_x)
=\gamma_V\varepsilon(1-\beta)
=\boxed{\varepsilon\sqrt{\frac{1-\beta}{1+\beta}}},
\qquad \beta=V/c.
\]
令 $\varepsilon'/\varepsilon=1/2$，
平方得
$(1-\beta)/(1+\beta)=1/4$，
故 $4-4\beta=1+\beta$，
\[
\boxed{V=\frac35c}.
\]
\excheck 与光子同向运动的观察者测得红移；
若 $V$ 逆光子方向，变换应得到蓝移而不是这一公式。

\exproblem{383}{61}
证明对于有静质量的粒子，组合
$E^2-p^2c^2$ 是所有惯性系中的不变量，
并给出这一不变量的数值。
\exsolution
在粒子静止系有
$E_0=m_0c^2$、$\bm p_0=\bm0$，
故
$E_0^2-p_0^2c^2=m_0^2c^4$。
由上一题的四动量洛伦兹变换，
\[
\begin{aligned}
E'^2-c^2p_x'^2
&=\gamma_V^2[(E-Vp_x)^2-
c^2(p_x-VE/c^2)^2]\\
&=\gamma_V^2(1-\beta^2)(E^2-c^2p_x^2)
=E^2-c^2p_x^2.
\end{aligned}
\]
垂直方向动量不变，因此一般三维情况下
$E'^2-c^2|\bm p'|^2=E^2-c^2|\bm p|^2$。
结论为
\[
\boxed{E^2-p^2c^2=m_0^2c^4}.
\]
\excheck 静止质量 $m_0$ 与惯性系无关，
而粒子的总能量、动量分别随参考系改变。

\exproblem{384}{61}
一颗静质量为 $m_0$ 的中子动能为
$T=2m_0c^2$，撞向另一颗实验室静止的中子。
求：(a) 两中子在质心系的总动能
$\widetilde T$ 与每颗中子的动量模
$\widetilde p$；
(b) 这个质心系相对实验室的速度。
\exsolution
实验室中，入射中子总能量
$E_1=m_0c^2+T=3m_0c^2$，
靶中子总能量 $E_2=m_0c^2$。
故系统总能量
$E_{\rm tot}=4m_0c^2$。
入射动量模由
$p_1c=\sqrt{E_1^2-m_0^2c^4}
=2\sqrt2\,m_0c^2$。
系统四动量不变量为
\[
E_{\rm cm}^2
=E_{\rm tot}^2-p_{\rm tot}^2c^2
=(16-8)m_0^2c^4=8m_0^2c^4.
\]
因此
$E_{\rm cm}=2\sqrt2\,m_0c^2$。
去掉两颗静止能量得质心系总动能：
\[
\boxed{\widetilde T
=2(\sqrt2-1)m_0c^2
\simeq777\,\mathrm{MeV}}.
\]
两中子在质心系总动量为零且质量相等，
每颗总能量为
$E_{\rm cm}/2=\sqrt2 m_0c^2$，
因此
\[
\boxed{\widetilde p=
\frac1c\sqrt{2m_0^2c^4-m_0^2c^4}
=m_0c\simeq940\,\mathrm{MeV}/c}.
\]
质心速度可由实验室总能量与总动量确定：
\[
\boxed{V_{\rm cm}
=\frac{p_{\rm tot}c^2}{E_{\rm tot}}
=\frac{c}{\sqrt2}
\simeq2.12\times10^8\,\mathrm{m/s}}.
\]
\excheck 必须使用系统的\emph{四动量不变量}
而非只在实验室中对两颗中子动能做算术平均。

\exproblem{385}{61}
静质量为 $m_0$、动能为 $T$ 的粒子
撞上一颗实验室静止、具有相同静质量的粒子，
两者完全结合为一个复合粒子。
求复合粒子的静质量 $M_0$
及它相对于实验室的速度 $V$。
\exsolution
碰前系统总能量与动量为
\[
E_{\rm tot}=(2m_0c^2+T),\qquad
p_{\rm tot}c=\sqrt{T(T+2m_0c^2)}.
\]
碰后复合粒子的静质量由四动量不变量定义：
\[
M_0^2c^4=E_{\rm tot}^2-p_{\rm tot}^2c^2
=(2m_0c^2+T)^2-T(T+2m_0c^2).
\]
展开后
$M_0^2c^4=4m_0^2c^4+2m_0c^2T$，
\[
\boxed{M_0=\sqrt{2m_0(2m_0+T/c^2)}}.
\]
复合粒子速度由
$\bm V=\bm p_{\rm tot}c^2/E_{\rm tot}$，
故
\[
\boxed{V=c\sqrt{\frac{T}{T+2m_0c^2}}}.
\]
\excheck $T>0$ 时 $M_0>2m_0$，
多出的静质量体现碰撞耗散的内部能量；
若把它直接写成 $2m_0$，违反总能量守恒。

\exproblem{386}{61}
问一个高能质子撞向实验室中\emph{静止}的另一质子时，
入射质子动能 $T_{\rm beam}$ 至少需要多高，
才能使该过程的\emph{质心系总动能}
等于两束相向质子、每束动能
$T=25.0\,\mathrm{GeV}$ 时的质心系总动能？
\exsolution
两质子相向且各自动能为 $T$ 时，
质心系就是两束的对称实验室系。
质心系总能量为
$E_{\rm cm}=2(m_pc^2+T)$。
对应的固定靶碰撞中，
实验室总能量
$E_{\rm lab}=2m_pc^2+T_{\rm beam}$，
入射动量满足
$p_{\rm lab}^2c^2
=T_{\rm beam}(T_{\rm beam}+2m_pc^2)$。
四动量不变量为
\[
s=E_{\rm lab}^2-p_{\rm lab}^2c^2
=4m_p^2c^4+2m_pc^2T_{\rm beam}.
\]
令它等于对撞束的
$s=4(m_pc^2+T)^2$，
得到
\[
2m_pc^2T_{\rm beam}
=4[(m_pc^2+T)^2-m_p^2c^4]
=4T(T+2m_pc^2).
\]
所以
\[
\boxed{T_{\rm beam}=
\frac{2T(T+2m_pc^2)}{m_pc^2}}.
\]
用质子静能 $m_pc^2\simeq0.938\,\mathrm{GeV}$：
\[
\boxed{T_{\rm beam}
\simeq\frac{2(25.0)(25.0+1.876)}{0.938}
\,\mathrm{GeV}\simeq1.43\times10^3\,\mathrm{GeV}
=1.43\,\mathrm{TeV}}.
\]
\excheck \textbf{本版答案单位冲突：}
书末公式与原题 $T=25\,\mathrm{GeV}$ 相容，
但数值旁印 $1.43\times10^3\,\mathrm{MeV}$；
按单位计算应为
$1.43\times10^3\,\mathrm{GeV}$，
相差1000倍。本稿保留物理量纲正确的值。

\exproblem{387}{61}
一颗静质量为 $m_0$ 的粒子
原来在实验室静止，
自发衰变为静质量分别为
$m_1,m_2,m_3$ 的三颗粒子。
求粒子1可能获得的最大总能量。
\exsolution
初态总四动量为
$(m_0c,\bm0)$。
将粒子2与3组成一个复合末态系统，
其不变质量记为 $M_{23}$。
总静质量下界
$M_{23}\ge m_2+m_3$；
仅当粒子2、3相对静止、
具有同一速度时取等。
把衰变视为
$m_0\to m_1+M_{23}$ 的二体末态。
在初始粒子静止系，由
$P_{23}=P_0-P_1$ 和
$P_{23}^2=M_{23}^2c^2$，
\[
M_{23}^2c^4
=(m_0c^2-E_1)^2-p_1^2c^2
=m_0^2c^4+m_1^2c^4-2m_0c^2E_1.
\]
于是
\[
E_1=
\frac{m_0^2+m_1^2-M_{23}^2}{2m_0}c^2.
\]
为使 $E_1$ 最大，令 $M_{23}=m_2+m_3$：
\[
\boxed{E_{1,\max}=
\frac{m_0^2+m_1^2-(m_2+m_3)^2}{2m_0}c^2}.
\]
\excheck 该值可达到的必要条件是
$m_0\ge m_1+m_2+m_3$；
要最大化单个粒子的能量，
\emph{其余两粒子须作为一个共动整体}
使不变质量最小。

\exproblem{388}{61}
一枚相对论火箭沿直线喷出气体，
气体相对火箭的速度大小固定为
$u$ 且满足 $u\ll c$。
火箭的静质量由初始 $m_0$
逐渐减到 $m$，无外力。
求火箭相对初始静止系的速度 $v(m)$。
\exsolution
每次考虑一个火箭的\emph{瞬时共动惯性系}。
以沿推进方向的火箭速度增量为正，
喷气速度为 $-u$。
在非相对论喷流近似下，
微元动量守恒得到
$m\,\mathrm dv_{\rm com}=-u\,\mathrm dm$，
其中 $\mathrm dm<0$。
为使不同瞬时共动系中的速度增量相加，
引入快度
$\chi=\operatorname{artanh}(v/c)$。
快度差与瞬时共动速度增量满足
$\mathrm d\chi=\mathrm dv_{\rm com}/c$，
因此
\[
\mathrm d\chi=-\frac{u}{c}\frac{\mathrm dm}{m}.
\]
从 $m_0$、初速零积分到 $m$：
\[
\chi=\frac uc\ln\frac{m_0}{m}.
\]
最后 $v/c=\tanh\chi$，
利用
$\tanh z=(1-e^{-2z})/(1+e^{-2z})$：
\[
\boxed{\frac vc=
\tanh\!\left(\frac uc\ln\frac{m_0}{m}\right)
=\frac{1-(m/m_0)^{2u/c}}
{1+(m/m_0)^{2u/c}}}.
\]
\excheck 当 $u/c$ 很小且累计速度也远小于 $c$，
$\tanh z\simeq z$，
退化为通常的火箭方程
$v\simeq u\ln(m_0/m)$；
在严格相对论模型中 $v<c$ 恒成立。
% === END 08-relativity-381-388.tex ===
